We implemented the FPPVS framework in Python. Our first version accepted only
about 10\% of requests, which is below random placement. The causes were an
unstated arrival rate that overloaded the network, a bunch of bugs, a sharing
model that could overbook capacity, and a loose reading of the paper (pruning by
compute only, any DC as a candidate, global pooling in the CNN). When we
implemented it closely from Section 4.3 of the paper, the original FPPVS works
much better. On top of that, we stated the workload as an explicit offered load
and added a look-ahead feasibility mask, an instance-overhead model for VNF
sharing, a per-data-center actor with a learnable placement prior,
rollback-and-retry admission and fixed federated training.

With these changes, updated FPPVS accepts \fppvsMinAcc--\fppvsMaxAcc\% of
requests on both topologies at heavy load (90--110\% of capacity), which is
above the 70\% target in every main scenario and every seed. On the same
scenarios, the original FPPVS comes second with \origMinAcc--\origMaxAcc\%,
ahead of the 2022 adaptation (\zhangMinAcc--\zhangMaxAcc\%) and of greedy and
random placement in every run. At full load, the original comes close to the
paper's published value on Metro (69.6\% against 75.9\%), but stays below it on
Backbone (72.9\% against 88.1\%). Coming to resources, updated FPPVS holds about
half the link bandwidth per running service compared to greedy and random, 11\%
less than the original FPPVS and about the same as the 2022 adaptation, and it
starts far fewer VNF instances than any other method. The ablation study shows
that VNF sharing and the problem modifications are what drive these gains, while
the learned part of the policy doesn't add much beyond its prior.

Future work: (1) evaluate on the real USAnet and Italian operator topologies and
traces once they are available; (2) re-implement the paper's baselines; (3) give
the agent a multi-request horizon, for example by treating a batch of arrivals
as one episode, so that learning can actually improve on the prior; and (4) add
VNF migration so that long-lived chains can be repacked when links fill up.
